% ============================================================
% !TeX root = ../main.tex
%  Chapter 3 — Methodology
% ============================================================
\chapter{Methodology}
\label{ch:methodology}

\noindent
This chapter presents the methodology for reconstructing three-dimensional LV
geometry from sparse SAX contours and measuring myocardial wall thickness on
the resulting surfaces. The approach combines synthetic geometric supervision
with clinical MRI-derived data, a phase-conditioned implicit representation,
and a comparison of local wall-thickness estimators.

\section{Overview}
\label{sec:method-overview}
\noindent
The input to the method is what a clinical short-axis (SAX) cardiac MRI
examination provides once the left ventricle (LV) has been outlined, or
segmented, on each image. In a SAX acquisition the heart is imaged in a stack of parallel
planes cut perpendicular to the long axis of the ventricle. On each of these
roughly ten planes the segmentation yields one closed ring of points on the
endocardium, the inner boundary of the blood-filled cavity, and one on the
epicardium, the outer boundary of the muscular wall, which is called the
myocardium. Each stack is labelled as acquired at end-diastole (ED), when the
ventricle is fullest, or at end-systole (ES), when it is most contracted.
Panels~(a) and~(b) of \cref{fig:methodology-pipeline} show such a stack. The
planes are several millimetres apart, so the wall is observed only where a
plane cuts it and has to be inferred everywhere else. The required outputs are
the two closed surfaces that pass through these rings and fill the gaps between
them (panel~(c)), a thickness value at every point of the reconstructed wall
(panel~(d)), and a summary of those values in the 17-segment layout of the
American Heart Association (AHA), the standard division of the ventricle used
in clinical reporting (panel~(e)).

\begin{figure}[H]
\centering
\includegraphics[width=0.85\linewidth]{cardiosdf_pipeline_visual_flow}
\vspace{2mm}
\caption[End-to-end reconstruction and thickness-measurement pipeline]{Overview
of the proposed approach, from sparse SAX observations to reconstructed LV
surfaces and wall-thickness outputs. Panels~(a)--(b) show ACDC data;
panels~(c)--(e) use illustrative SSM-derived geometry with a synthetic
epicardium and do not show reported results.}
\label{fig:methodology-pipeline}
\end{figure}

Three obstacles shape the methodology. The first is that no public dataset
pairs such contour stacks with independently acquired three-dimensional LV
surfaces. The examples the model needs to learn from do not exist and had to
be constructed. Three-dimensional LV surfaces were therefore built from two
sources. The first is a statistical shape model (SSM), a compact description
of how ventricle shape varies across a population. Sampling it gives complete
surfaces, which are used to pre-train the model, that is, to teach it the
general shape of a ventricle before it sees any patient. The second source is
the clinical contours of two public datasets, ACDC and M\&Ms-2, through which
smooth surfaces are fitted. These give patient-specific ED and ES targets for
fine-tuning, the second training stage, in which the model adapts to real
patients.

The second obstacle is that the two surfaces are not independent: the
epicardium must lie outside the endocardium everywhere, including between the
slices where nothing is observed. The model therefore represents the
endocardium by a signed-distance field (SDF), a function that gives at every
point in space the distance to the surface, negative inside the cavity and
positive outside. The epicardial field is obtained by subtracting a strictly
positive offset from it, so the two surfaces cannot cross. The model also
receives the cardiac phase as an input, so a single model can represent ED and
ES.

The third obstacle is that wall thickness has no unique definition on a curved
wall: the value depends on how a point on the inner surface is paired with a
point on the outer surface. Instead of fixing one definition in advance, four
established estimators are run on the same reconstructed surfaces, the most
consistent one is selected, and only then are regional AHA-17 values computed.

Most of the building blocks are established methods, and the chapter names
their origin where each appears. The model has two parts: an encoder, which
reads the contour points and condenses them into a short description, and a
decoder, which turns that description into the two fields. The decoder is of
the DeepSDF type \parencite{deepsdf2019}: given the description of a shape, it
returns the signed distance at any point, so the surface can be read out at
any resolution. The encoder is a PointNet \parencite{qi2017pointnet}: it
applies the same small network to every point and keeps, for each learned
feature, the strongest response, so the result does not depend on how many
points there are or in what order they arrive. Because such a summary forgets
where each point was, a second encoder keeps a coarse three-dimensional grid
of local features, as in Convolutional Occupancy Networks
\parencite{peng2020convonet}. Two standard ingredients complete the model.
Fourier features \parencite{tancik2020fourierfeatures} present the coordinates
of a point as sine and cosine waves, which makes fine detail easier to learn,
and the regularisation terms of \textcite{gropp2020igr} make the predicted
field behave like a true distance. Outside the model, the clinical targets are
built with the radial basis function (RBF) surface fitting of
\textcite{carr2001rbf}, a classical way to pass a smooth surface through
scattered points; the population shape model is the UK Digital Heart Project
atlas of \textcite{bai2015biventricularatlas}; and the thickness estimators
and the AHA-17 layout come from the literature reviewed in
\cref{subsec:thickness-review}.

What is new is how these parts are put together to close the gap identified
in \cref{subsec:research-gap}. Every contour point carries its tissue and its
cardiac phase, so one model serves ED and ES. Endocardial and epicardial
points are summarised separately before they are merged. The epicardial field
is tied to the endocardial field by a positive offset, which guarantees a
positive wall and is trained directly against the measured wall distance.
Training moves from synthetic SSM geometry to clinical ED and ES targets
through one shared sample format (cache). And the evaluation is designed for a setting
in which no three-dimensional ground truth exists.
\Cref{tab:model-components-origin} in \cref{sec:architecture} summarises this
division for the model itself.

The chapter follows the pipeline of \cref{fig:methodology-pipeline}.
\Cref{sec:data-preparation} constructs the two supervision streams and the
shared sample format. \Cref{sec:architecture} describes the model, and
\cref{sec:training-inference} how it is trained, how surfaces are extracted
from it, and which alternative reconstructions it is compared against.
\Cref{sec:wall-thickness-protocol} defines the thickness estimators, the
three-dimensional thickness maps, and the AHA-17 aggregation. \Cref{ch:results}
then applies this protocol: it measures how closely each reconstruction follows
the observed rings and how far it agrees with an RBF fit and an SSM fit of the
same rings, how the reconstructions degrade as slices are removed, how the
four thickness estimators agree with one another, and how regional thickness
differs between healthy hearts and hearts with hypertrophic cardiomyopathy, a
disease in which the wall thickens abnormally, and between ED and ES. These
analyses answer, in turn, the three research questions of
\cref{sec:research-questions}.

\section{Data Preparation}
\label{sec:data-preparation}
\noindent
The model is trained to map contours to surfaces, so every training case must
contain both: a sparse contour stack as input and a pair of closed
three-dimensional surfaces as target. Public cardiac MRI datasets provide only
the first half. ACDC and M\&Ms-2 supply images and slice-wise segmentation
masks, label images in which every pixel is marked as cavity, myocardium, or
background, but not the closed (``watertight'') endocardial and epicardial
surfaces that a signed-distance decoder needs as supervision. The training data were
therefore constructed from two sources that fill this gap in complementary
ways. The UK Digital Heart Project LV
SSM\footnote{\url{https://github.com/UK-Digital-Heart-Project/Statistical-Shape-Model}}
supplies complete surfaces over a controlled range of healthy shapes, from
which 1300 synthetic ED cases were generated for pre-training. The clinical
segmentations supply patient-specific anatomy in both cardiac phases, and were
converted into closed surfaces by fitting smooth implicit functions to their
contours, for fine-tuning. The two sources are described in turn, followed by
the shared format, augmentation, and patient-level partitions through which
they enter training.

\subsection{Synthetic ED geometry from the statistical shape model}
\label{subsec:missing-ground-truth}
\noindent
A statistical shape model describes a population of shapes by a mean shape and
a small number of principal modes of variation, the main independent ways in
which the shapes of the population differ from that mean
\parencite{heimann2009ssmreview}.
Sampling it yields complete, closed surfaces whose variation is known and
controlled, which makes it a convenient source of supervision for pre-training.
What it cannot provide matters just as much: the sampled shapes carry no
scanner-dependent acquisition effects, no variation in how a clinician would
have drawn the contours, and, for the model used here, no systolic geometry.
The synthetic stream is therefore used to teach the model what an LV looks
like in general, before the clinical stream exposes it to the conditions of
the evaluation data.

The shape model is the UK Digital Heart Project LV SSM of
\textcite{bai2015biventricularatlas}, built from more than 1,000
high-resolution cardiac MR images of healthy subjects. It provides a mean
endocardial surface of 22\,043 vertices (the corner points of the triangles
that make up the mesh) together with 100 principal modes of variation, stored
as a $66\,129\times100$ matrix of mode vectors, and a corresponding vector of
100 eigenvalues $\lambda_i$, each measuring how much of the population's
variation its mode explains. A new endocardial mesh is generated by
\begin{equation}
	X(b) = \bar{X} + \Phi b,
	\label{eq:ssm-generation}
\end{equation}
where $\bar{X}$ is the mean shape flattened to a vector, $\Phi$ is the matrix of
mode vectors, and $b = [b_1, \ldots, b_{100}]^\top$ is the sampled coefficient
vector. The standard deviation $\sigma_i = \sqrt{\lambda_i}$ sets the scale of
variation along mode $i$. The first modes capture the largest sources of
variation in the population: mode~1 mainly controls overall LV size, mode~2 the
elongation along the long axis, and higher modes increasingly local
deformations (\cref{fig:ssm-mean-modes}).

\begin{figure}[H]
	\centering
	\includegraphics[width=0.9\linewidth]{fig_ssm_mean_modes}
	\vspace{2mm}
	\caption[LV SSM mean shape and principal modes]{UK Digital Heart Project LV
	SSM mean shape and first three PCA
	modes at $\pm3\sigma$.}
	\label{fig:ssm-mean-modes}
\end{figure}

Of the 100 modes, the first 32 capture 99\% of the shape variance and define
the sampling space. Each candidate coefficient vector $b$ was drawn by sampling
each active component as $b_i \sim \mathcal{N}(0, \sigma_i^2)$ and was accepted
only if it satisfied
\begin{equation}
	\left|\frac{b_i}{\sigma_i}\right| \leq 3,
	\qquad
	\sum_i \left(\frac{b_i}{\sigma_i}\right)^2 \leq \chi^2_{0.99,d},
	\label{eq:ssm-latent-constraint}
\end{equation}
where $d$ is the number of active modes and $\chi^2_{0.99,d}$ is the 99th
percentile of the chi-squared distribution with $d$ degrees of freedom. The
first condition keeps every individual mode within three standard deviations of
the mean. The second bounds the overall deviation: the sum measures how far the
sample lies from the mean shape once every mode is weighed by its own spread
(the Mahalanobis distance), and the threshold admits only combinations that
would occur in 99\% of the population. Each accepted vector is substituted into
\cref{eq:ssm-generation} to obtain new vertex positions, which are connected
into triangles exactly as in the mean mesh; the result is a closed and
consistently oriented surface by construction. Each surface is then checked
against four simple geometric descriptors, namely its enclosed volume, its
surface area, a sphericity index (how close the shape is to a sphere), and the
direction of its surface normals (the directions perpendicular to the surface)
at the apex and base, and it is discarded if any value falls outside bounds
calibrated on the population; exact volume and area computations on every
fortieth sample confirmed that the fast descriptors were reliable. Of 2048
candidates, 1300 (63\%) were accepted.

The SSM provides only the endocardial surface. An epicardial surface was
synthesised by moving each endocardial vertex $p_i$ outward along its normal
$n_i$ by a distance $d_i$,
\begin{equation}
	q_i = p_i + d_i n_i,
	\label{eq:synthetic-epi-offset}
\end{equation}
with $d_i$ set to approximately 10~mm near the base and 5~mm near the apex,
plus Gaussian noise with a standard deviation of 0.5~mm so that the wall is not
perfectly uniform. This synthetic epicardium exists only to give the decoder a
consistent outer surface during pre-training, so that the model learns the
general relation between an inner and an outer LV surface before it sees
clinical data. It does not reproduce clinical wall thickness and is never used
as a reference for any thickness result.

Finally, each surface pair was converted into the form the model will see.
Slicing both surfaces at 10 evenly spaced levels along the long axis produces
sparse SAX contour rings that mimic a clinical acquisition. For each case, 2048
query points, locations in space at which the decoder will be asked for a
value during training, were sampled near the surfaces and in the surrounding
volume, and the signed distance to each surface, or an inside/outside label,
was recorded at every one of them; these points provide the dense volumetric
supervision for the decoder. \Cref{fig:ssm-sampling-pipeline} summarises the whole procedure, from
the accepted region of coefficient space to the contour rings and query points
of one case. The synthetic stream fixes the general shape of the problem; the
clinical stream, described next, supplies the anatomy on which the model is
evaluated.

\begin{figure}[H]
	\centering
	\includegraphics[width=0.9\linewidth]{fig_ssm_sampling_pipeline}
	\vspace{2mm}
	\caption[Synthetic ED dataset generation from SSM sampling.]{Synthetic ED
	dataset generation, from SSM sampling to geometric supervision. Top row,
	statistical constraints: (a)~distribution of the overall-deviation
	(Mahalanobis) statistic of \cref{eq:ssm-latent-constraint} with the
	$\chi^2_{0.99,32}$ acceptance threshold; (b)~cumulative shape variance, with
	32 active modes reaching 99\%; (c)~accepted coefficient samples in the first
	two modes; (d)~distribution of the synthetic epicardial offset of
	\cref{eq:synthetic-epi-offset}. Bottom row, geometric construction: (e)~mesh
	synthesis via \cref{eq:ssm-generation}; (f)~an accepted shape next to a
	rejected outlier; (g)~the 10 SAX contour levels; (h)~the 2048 near-surface
	and volumetric query points used for decoder supervision.}
	\label{fig:ssm-sampling-pipeline}
\end{figure}

\subsection{Clinical ED and ES surfaces from SAX labels}
\label{subsec:real-ed-es-datasets}
\noindent
The clinical data come from two public collections, the Automated Cardiac
Diagnosis Challenge (ACDC) \parencite{bernard2018acdc} and the second
Multi-Centre, Multi-Vendor and Multi-Disease Cardiac Segmentation challenge
(M\&Ms-2) \parencite{martinisla2023mnms2}, which together provide 460 patients
(100 from ACDC and 360 from M\&Ms-2) with outlines of the cavity, the
myocardium, and the right ventricle drawn by experts at ED and ES.
Both are available for research at no cost after registration and acceptance
of a data-use agreement, and the shape model of the previous section is
likewise freely available. These data contribute what the synthetic stream
lacks: scanner and site variability, the way clinicians actually delineate the
wall, patient-specific and in part pathological anatomy, and the contracted
geometry of end-systole.

Segmentation masks are not yet surfaces. \Cref{fig:problem-visualisation}
shows representative ACDC slices with the cavity, myocardium, and
right-ventricular labels, and the ED labels of one patient stacked at their
true through-plane positions. The labels are defined one slice at a time, and
the slices are often several millimetres apart, so stacking them gives a set of
sparse rings with nothing in between. Turning these rings into closed surfaces
is a reconstruction problem in itself, and the method chosen for it determines
what the model is trained to reproduce.

\begin{figure}[H]
	\centering
	\includegraphics[width=0.98\linewidth]{acdc_problem_visualisation}
	\vspace{2mm}
	\caption[Clinical input for the real-data reconstruction pipeline.]{Clinical-input view used in the real-data pipeline: (a) representative
	ACDC ED/ES SAX slices with labels and (b) the corresponding sparse ED SAX
	stack in physical three-dimensional spacing.}
	\label{fig:problem-visualisation}
\end{figure}

Three options were available. Connecting corresponding points on consecutive
rings (lofting) requires a one-to-one pairing of vertices between rings that
the data do not provide, and gives a surface that is only piecewise linear
between slices. Fitting the SSM would impose a population prior on every
target and, because the available model has ED modes only, would either force
systolic hearts into a diastolic shape space or treat the two phases
differently. Fitting a smooth implicit function to the rings has neither
drawback: it needs no vertex pairing and no atlas, and the same procedure
applies to ED and ES alike. The clinical targets were therefore built by
Radial Basis Function (RBF) fitting, which leaves the SSM in its role as the
source of synthetic pre-training geometry.

Each segmentation is distributed as a NIfTI file, the standard container for
medical image volumes. The file stores the label volume together with the
physical size and position of its voxels, the three-dimensional pixels of the
volume. The volume is first reoriented to a common anatomical convention,
called RAS+, in which the three axes point towards the patient's right,
anterior, and superior sides; the physical voxel spacing is preserved, so that
every patient is expressed in the same frame and in millimetres. Each voxel of
the label volume carries an integer that names the structure it belongs to,
such as cavity, myocardium, or right ventricle; because the two datasets
number these structures differently, the labels are mapped to a common
definition of cavity and myocardium. On every acquired SAX slice, the boundary
of the cavity gives the endocardial ring and the outer boundary of cavity plus
myocardium gives the epicardial ring; when a slice contains more than one
closed boundary, for instance because of a small detached fragment of the
label, only the largest ring is kept. Each ring is resampled to 60 evenly
spaced points, so that every slice contributes the same number of points
regardless of image resolution, and the points are stored at the measured
position of the slice along the long axis. These points serve two purposes:
they are the sparse observation the encoder receives, and they are the
constraints from which the target surfaces are built.

The endocardial and epicardial surfaces are fitted separately. For either
surface $s$, let $p_i$ be a contour point and $n_i$ its outward in-plane
normal, estimated from the local contour tangent and pointing away from the
ring centre. Following \textcite{carr2001rbf}, a smooth function $f_s$ is
defined over the whole volume as a sum of radial basis functions (RBFs),
smooth bumps whose value depends only on the distance to a centre $c_j$, plus
a linear term, and the surface is taken as its zero level set, the set of
points where the function is zero. With the thin-plate-spline kernel $\phi$
used here, the function is
\begin{equation}
	f_s(x) = \sum_{j=1}^{K} w_j\phi(\lVert x-c_j\rVert_2)
	+ a_0 + a^\top x,
	\qquad \phi(r)=r^2\log r,
	\label{eq:clinical-rbf-field}
\end{equation}
where the centres $c_j$ are the constraint points and the weights $w_j$,
$a_0$, and $a$ are chosen so that the function satisfies the constraints. The
contour points fix where the function is zero, and paired points placed a
distance $d=2.5$~mm inside and outside each contour point fix its sign and
local scale:
\begin{equation}
	f_s(p_i)=0,
	\qquad f_s(p_i+d n_i)=d,
	\qquad f_s(p_i-d n_i)=-d.
	\label{eq:clinical-rbf-constraints}
\end{equation}
A regularisation parameter of 0.5 lets the function pass close to, rather than
exactly through, the constraints, which reduces its sensitivity to small
irregularities in the drawn contours.

Each fitted function is evaluated on a 1.0~mm grid that extends 8.0~mm beyond
the rings in-plane and 1.5~mm beyond the first and last ring along the long
axis, so that the surface can close just beyond the observed extent. The region
where $f_s\leq0$, the inside of the surface, is reduced to its valid connected
part (small disconnected pockets are discarded) and converted to a
signed-distance volume, and the volume is smoothed with
a Gaussian filter of width 1.2~mm, which replaces each value by a weighted
average of its neighbours, before marching cubes, a standard algorithm that
extracts a triangle mesh from the zero level set on a grid, produces the
surface \parencite{lorensen1987marchingcubes}. Smoothing the field at this
stage removes the small step-like ripples that the voxel grid would otherwise
imprint on the surface, without moving mesh vertices afterwards. The extracted
mesh is accepted unchanged if it is a single closed shell with consistently
oriented faces and no handles or tunnels; otherwise only the failed property
is repaired. Cases with fewer than two valid rings per surface, a failed fit,
a surface that could not be repaired into a single closed shell, implausible
dimensions, or an endocardium not fully enclosed by its epicardium are
excluded. The full sequence of operations and their parameters is listed in
\cref{tab:real-pipeline} of \cref{app:supplementary-methods}, and
\cref{fig:real-mesh-steps} illustrates it on one case: the endocardial and
epicardial constraints are fitted separately, and their surfaces then pass
through the same checks.

\begin{figure}[H]
	\centering
	\includegraphics[width=0.98\linewidth]{fig_real_mesh_steps}
	\vspace{2mm}
	\caption[Clinical RBF surface-construction pipeline]{Clinical RBF
	contour-to-surface pipeline. Panel~(a) shows one acquired SAX slice with the expert
	endocardial, epicardial, and right-ventricular contours. Panel~(b) shows the
	corresponding basal, mid-cavity, and apical labels, while panels~(c) and
	(c$'$) stack the extracted contour rings.
	Panels~(d) and~(d$'$) overlay those constraints on the zero-level surfaces
	extracted from the smoothed RBF fields, and panels~(e) and~(e$'$) show the
	final topology-checked surfaces used to construct the clinical supervision.}
	\label{fig:real-mesh-steps}
\end{figure}

The two meshes of each case are centred and scaled with the same translation
and scale factor, and 2048 query points are drawn near the surfaces and
throughout a box that encloses them with a margin, with their signed distances
recorded, exactly as for the synthetic stream.

One property of these targets governs the interpretation of every later
result. They are \emph{surrogate targets}: RBF-derived surfaces that stand in
for the three-dimensional reference geometry that does not exist, and whose
shape between the slices depends on the kernel, the off-surface constraints,
the regularisation, and the extraction resolution just described. Agreement
between the model and these targets shows that the model has learned to
reproduce a particular, reasonable reconstruction of the rings, not that it has
recovered the true anatomy between them. \Cref{ch:results} keeps this
distinction explicit by also measuring every reconstruction against the
observed rings themselves.

\subsection{Shared sample format, augmentation, and splits}
\label{subsec:shared-dataset-assembly}
\noindent
Both streams enter training through the same sample format, the same
augmentation rule, and partitions defined at the level of patients.

Every case, synthetic or clinical, is stored as one item with the fields listed
in \cref{tab:cache-schema}: the contour points with their tissue and phase
labels, surface samples and query points derived from the target meshes, and
the coordinate information that keeps observation and target in the same
frame. The ED and ES items contain 10 SAX levels, 60 points per available ring,
and 2048 query points per case. The format deliberately separates what the
model observes from what it learns from. The encoder receives only the first
three fields; the decoder is supervised by the remaining ones, which are never
shown to the encoder. At inference time the model therefore needs nothing but
the contours, and the architecture and loss described in the next two sections
do not change when training moves from the synthetic to the clinical stream.

\begin{table}[H]
	\centering
	\caption[Shared sample format]{Fields of one stored sample, shared by the
	synthetic and clinical streams, and their role in training.}
	\label{tab:cache-schema}
	\small
	\renewcommand{\arraystretch}{1.12}
	\begin{tabular}{p{0.20\textwidth} p{0.25\textwidth} p{0.43\textwidth}}
		\toprule
		Field & Origin & Role \\
		\midrule
		Contour points & Synthetic mesh slicing or clinical slice labels & Sparse SAX observation for the encoder \\
		Tissue label & Contour extraction & Distinguishes endocardial from epicardial points during tissue-specific pooling \\
		Phase label & Dataset metadata & Conditions the encoder on ED or ES geometry \\
		Surface samples & Prepared endocardial and epicardial meshes & Surface-fitting and wall-offset supervision \\
		Query points and targets & Near-surface and volumetric mesh sampling & Dense signed-distance and off-surface supervision \\
		Mesh metadata & Preprocessing pipeline & Coordinates, scale, faces, and orientation for evaluation and visualisation \\
		\bottomrule
	\end{tabular}
\end{table}

The clinical stream is small, and its contours vary in how they were drawn and
acquired. To make the encoder robust to such variation, the observed contours
are perturbed during training while the targets are left untouched. For each
training sample, the data loader applies small independent in-plane shifts to
each slice, moving the endocardial and epicardial rings of a slice together so
that their relation is preserved, as well as small random displacements of
individual points, a global change of scale, and random removal of contour
points; with lower probability it also drops whole slices, always keeping at
least three, or rotates the stack about the long axis. The target mesh, surface
samples, query points, and tissue and phase labels stay fixed, so the encoder
must recover one fixed target from
many plausible versions of its sparse observation
(\cref{fig:contour-augmentation}). Such label-preserving perturbations are a
standard way to improve robustness when labelled data are limited
\parencite{shorten2019augmentation,chlap2021augmentation}. Augmentation is
applied to both streams and only after the data have been split, so that no
augmented view of a test patient is ever seen in training.

\begin{figure}[H]
	\centering
	\includegraphics[width=0.9\linewidth]{fig_contour_augmentation}
	\vspace{2mm}
	\caption[Contour-input augmentation examples]{Examples of contour-input
	augmentation: the unaltered contour stack, slice translation, slice
	rotation, and contour-point dropout. Endocardial and epicardial points are
	shown in blue and orange, respectively; the target mesh is the same in
	every case.}
	\label{fig:contour-augmentation}
\end{figure}

Not every clinical patient yields a usable target. A study may contain only
long-axis views or too few slices, or its segmentation may fail the
plausibility checks of \cref{subsec:real-ed-es-datasets}. After these checks,
444 of the 460 patients provide a valid ED mesh and 423 a valid ES mesh;
together with the 1300 accepted synthetic meshes they form the three streams
of \cref{tab:dataset-summary}. The clinical patients are divided into
training, validation, and test sets in a 70\,\%/15\,\%/15\,\% proportion,
rounded to whole patients and applied to patients rather than to meshes, so
that the ED and ES meshes of one patient always fall in the same set and no
information passes from the test patients into training. The synthetic meshes
are used for pre-training only and never appear in the validation or test
sets, so every number reported in \cref{ch:results} is computed on clinical
cases. The role of each stream in the two training stages is summarised in
\cref{app:supplementary-methods}. With every case in one format, the model can
now be described without reference to the source of a sample.

\begin{table}[H]
	\centering
	\caption[Dataset composition and patient-level splits.]{Final sample counts and partitioning of the three data
		streams. Real streams report patients that pass the quality checks; the
		synthetic stream reports meshes accepted out of the 2048 sampled shape
		vectors. The reported clinical evaluation contains no synthetic cases.}
	\label{tab:dataset-summary}
	\small
	\begin{tabular}{l l c c}
		\toprule
			Stream & Source & Phase & Count \\
		\midrule
			Real data & ACDC and M\&Ms-2 & ED & 444 patients \\
			Real data & ACDC and M\&Ms-2 & ES & 423 patients \\
			Synthetic data & UK Digital Heart SSM & ED & 1300 meshes \\
		\bottomrule
	\end{tabular}
\end{table}

\section{Model Architecture}
\label{sec:architecture}
\noindent
The data impose three requirements on the model. Its input is a set of a few
hundred points whose number varies from case to case and whose order carries
no meaning, so the encoder must be indifferent to point count and order. Its
output is a pair of surfaces that must be continuous between the slices and
correctly nested everywhere, so the decoder must produce two related fields
rather than one. And the same model must represent diastolic and systolic
hearts, so the phase must reach every stage that interprets the contours.
\Cref{fig:model-architecture} shows the resulting design. A global encoder
summarises the whole contour set in one vector, a local encoder records where
in space the contours lie, and an implicit decoder, evaluated at any point,
combines both with the coordinates of that point to return the endocardial
signed distance and the positive offset that defines the epicardial field.

\begin{figure}[H]
	\centering
	\includegraphics[width=0.78\linewidth]{fig_model_architecture}
	\vspace{2mm}
	\caption[Proposed model architecture.]{Phase-conditioned model architecture,
	from the sparse contour observation~(a) through the global and local encoder
	paths~(b) and the implicit decoder~(c) to the coupled fields and extracted
	surfaces~(d). The point cloud in~(a) and the cut surfaces in~(d) are the
	observed contours and the corresponding reconstruction of one ACDC case,
	with the epicardium opened to expose the endocardial cavity. The offset
	$\delta$ shown in the figure is a separation between field values, not a
	measured thickness.}
	\label{fig:model-architecture}
\end{figure}

Each of these parts has a precedent, and \cref{tab:model-components-origin}
separates what is borrowed from what is specific to this work. The decoder is a
signed-distance network conditioned on a latent code, a compact vector of
numbers that summarises one shape, as introduced by DeepSDF
\parencite{deepsdf2019}; here, however, the code is produced by an encoder from
the observed contours rather than fitted separately to every shape by
optimisation, so that a new case is reconstructed in a single evaluation of the
network. The global encoder is a PointNet \parencite{qi2017pointnet}, in which
the same small network is applied to every point and the results are combined
by max-pooling, keeping for each feature the largest value over all points;
this is the standard way to encode a set of points whose order carries no
meaning. The local feature volume follows Convolutional Occupancy Networks
\parencite{peng2020convonet}, which showed that a single summary vector loses
the information of where each piece of evidence was, and that spreading the
point features over a coarse three-dimensional grid and reading that grid near
each query point restores it. The positional encoding is the
Fourier-feature mapping of \textcite{tancik2020fourierfeatures}. Compared with
the recent cardiac implicit models reviewed in
\cref{subsec:dl-reconstruction-review}, which reconstruct from image data
\parencite{bras2026anisotropy} or from sparse segmentations of both ventricles
with long-axis views and a learned biventricular prior
\parencite{muffoletto2025nihc}, the proposed model takes only LV contour points
with a tissue and a phase label, and it makes the relation between the two
surfaces part of the architecture: the epicardial field is defined from the
endocardial field through a positive offset that is supervised directly.

\begin{table}[H]
	\centering
	\caption[Origin of the model components]{Model components, the prior work
	each one builds on, and what is specific to this work.}
	\label{tab:model-components-origin}
	\small
	\renewcommand{\arraystretch}{1.12}
	\begin{tabular}{p{0.19\textwidth} p{0.31\textwidth} p{0.40\textwidth}}
		\toprule
		Component & Builds on & Specific to this work \\
		\midrule
		Input representation & Unordered point-set input \parencite{qi2017pointnet} & Every point carries a tissue label and the cardiac phase, so one model serves ED and ES \\
		Global encoder & Shared per-point network with max-pooling \parencite{qi2017pointnet} & Endocardial and epicardial points are pooled separately before being merged into the latent code \\
		Local feature volume & Convolutional feature grid read at the query point \parencite{peng2020convonet} & Built from a few hundred contour points averaged into a coarse $16^3$ grid; fed to the decoder trunk and to the offset head \\
		Implicit decoder & Latent-conditioned signed-distance network \parencite{deepsdf2019} with Fourier features \parencite{tancik2020fourierfeatures} & Latent code supplied by the contour encoder instead of per-shape optimisation \\
		Output coupling & Surfaces as zero level sets of signed-distance fields \parencite{deepsdf2019} & Epicardial field defined as the endocardial field minus a strictly positive offset $\delta$, which a wall-distance loss supervises directly \\
		Training & Surface and distance fitting \parencite{deepsdf2019}; Eikonal regularisation \parencite{gropp2020igr} & Two stages, synthetic SSM pre-training then clinical ED/ES fine-tuning, through one shared sample format \\
		\bottomrule
	\end{tabular}
\end{table}

\subsection{What the model observes}
\label{subsec:architecture-input}
\noindent
For a batch of $B$ cases processed together, the encoder receives the contour
points as one padded array, or tensor,
\begin{equation}
	P \in \mathbb{R}^{B\times N\times5}
	\label{eq:contour-input-tensor}
\end{equation}
and a validity mask $M\in\{0,1\}^{B\times N}$ that marks which of the $N$ rows
are real points; padding only allows cases with different numbers of points to
share a batch. Each valid row is $p_i=(x_i,y_i,z_i,t_i,q_i)$: normalised
spatial coordinates, a tissue label $t_i$ that distinguishes endocardial from
epicardial points, and a phase label $q_i$ that identifies ED or ES. The
network receives neither image intensities nor a template mesh; everything it
knows about a case comes from these five values per point. Attaching the phase
to every point, rather than introducing it only at the decoder, lets both
encoders interpret the same arrangement of contours differently for a
diastolic and a systolic heart. The tissue label serves a different purpose: it
allows evidence about the cavity boundary and evidence about the outer boundary
to be summarised separately before they are combined.

\subsection{Contour encoder}
\label{subsec:contour-encoder}
\noindent
The global pathway is a PointNet \parencite{qi2017pointnet}: a shared network
with layer widths $5\rightarrow64\rightarrow128\rightarrow256\rightarrow256$,
with ReLU activations (simple non-linear functions) between the layers,
processes every point independently, so that no ordering is imposed within a
ring or between slices, and the per-point features are then pooled. Writing
$P = \{p_i\}_{i=1}^{N}$ for the points of one case, each point receives the
256-dimensional feature
\begin{equation}
	h_i = \phi_{\mathrm{enc}}(p_i),
	\label{eq:point-encoding}
\end{equation}
where $\phi_{\mathrm{enc}}$ is the shared network. The departure from the
standard PointNet is in the pooling. Instead of one maximum over all points,
the maximum is taken separately over the endocardial set $\mathcal{E}$ and the
epicardial set $\mathcal{P}$,
\begin{equation}
	z_{\mathrm{endo}} = \max_{i \in \mathcal{E}} h_i,
	\qquad
	z_{\mathrm{epi}} = \max_{i \in \mathcal{P}} h_i,
	\label{eq:tissue-pooling}
\end{equation}
independently in every feature channel, so that the strongest response to each
learned geometric pattern is kept for each boundary before the two are merged.
The two descriptors are concatenated into a 512-dimensional vector and
projected to the global latent code
\begin{equation}
	z = W_z [z_{\mathrm{endo}}, z_{\mathrm{epi}}] + b_z,
	\label{eq:latent-projection}
\end{equation}
with $z\in\mathbb{R}^{256}$, which gives the decoder a compact description of
the whole case.

Max-pooling has a cost: compressing all points into one vector discards where
each piece of evidence was located. The local pathway restores this
information, following Convolutional Occupancy Networks
\parencite{peng2020convonet}. A second shared network with widths
$5\rightarrow64\rightarrow64\rightarrow32$ assigns a 32-dimensional feature to
every valid point, and the features are scattered into a regular $16^3$ grid
covering the normalised extent $[-1.8,1.8]^3$; the features of points that fall
in the same cell are averaged, so the content of a cell does not depend on how
many contour points it happens to contain. Three $3\times3\times3$
convolutions, small filters that combine each cell with its immediate
neighbours, with Group Normalisation (a rescaling of the features that
stabilises training) and ReLU after the first two, spread information into
neighbouring cells and produce
\begin{equation}
	V \in \mathbb{R}^{32\times16\times16\times16}.
	\label{eq:local-feature-volume}
\end{equation}
For a query point $x$, the local feature is read from this volume by trilinear
interpolation, a weighted average of the eight surrounding cells,
\begin{equation}
	v(x) = \mathrm{Interp}(V, x).
	\label{eq:local-volume-interp}
\end{equation}
Unlike $z$, which is the same for every query of a case, $v(x)\in\mathbb{R}^{32}$
changes with the query location and describes the contour evidence near it.
The grid is deliberately coarse: a few hundred contour points cannot support a
fine volume, and its role is to give context to the decoder, not to serve as
the reconstruction grid.

\subsection{Implicit decoder with positive-offset coupling}
\label{subsec:monotone-epi-decoder}
\noindent
The decoder represents each surface as the zero level set of a function of
continuous coordinates, in the manner of DeepSDF \parencite{deepsdf2019}. A
network fed with plain coordinates tends towards overly smooth fields, so each
query point is first expanded with Fourier features
\parencite{tancik2020fourierfeatures}, which add sine and cosine components at
several spatial frequencies. For $x\in\mathbb{R}^3$, the encoding is
\begin{equation}
	\gamma(x) = [x, \sin(2^0\pi x), \cos(2^0\pi x), \ldots,
	\sin(2^{L-1}\pi x), \cos(2^{L-1}\pi x)],
	\label{eq:fourier-encoding}
\end{equation}
with $L=3$, giving $3+6L=21$ components. The encoded coordinate is concatenated
with the global code to form the decoder input
\begin{equation}
	h_{\mathrm{in}}(x)=[z,\gamma(x)]\in\mathbb{R}^{277}.
	\label{eq:decoder-input}
\end{equation}
An input layer maps this vector to width 512, and the main body of the decoder,
its trunk, applies eight fully connected layers of the same width with Softplus
activations, smooth non-linear functions that keep the field smooth and
differentiable. At the fourth hidden layer, the 277-dimensional input is
appended to the current state and projected back to width 512; this skip
connection lets the deeper layers see the case code and the coordinate
directly rather than through the whole trunk. The local feature $v(x)$ enters
the decoder at the input layer, at this skip connection, and in the offset
head described next, so that the prediction at a point depends on the nearby
contours as well as on the global shape.

The output stage is where the model departs from an ordinary implicit decoder.
Predicting the endocardium and the epicardium as two unrelated fields would let
them cross or touch wherever the data are sparse, producing a wall of zero or
negative thickness between slices. The final hidden representation therefore
feeds two coupled output layers, or heads: one predicts the endocardial
signed-distance value $f_{\mathrm{endo}}(x)$, the other a strictly positive
offset $\delta(x)$, and the epicardial field is defined as
\begin{equation}
	f_{\mathrm{epi}}(x) = f_{\mathrm{endo}}(x) - \delta(x),
	\qquad \delta(x)>0.
	\label{eq:monotone-epi}
\end{equation}
Positivity is enforced in the implementation by a sigmoid-based mapping, which
squashes any real number into a bounded positive range, with additional
headroom for large offsets. Because the epicardial field is always
lower than the endocardial field by $\delta$, its zero level set always lies
outside the endocardial one, which rules out wall collapse by construction. Two
qualifications apply. The constraint does not by itself guarantee an
anatomically valid surface or a valid mesh once the field has been sampled on a
grid, which is why the extracted meshes are checked and repaired in
\cref{subsec:inference-offset-extraction}. And $\delta$ is a difference between
two field values at one point, not the distance between the two extracted
surfaces, so it is not a measurement of wall thickness; thickness is measured
on the extracted surfaces with the methods of
\cref{sec:wall-thickness-protocol}. How the two fields are turned into meshes
is described in \cref{subsec:inference-offset-extraction}; the complete
configuration of the model is listed in \cref{tab:model-hyperparameters} of
\cref{app:supplementary-methods}.

\section{Training and Inference}
\label{sec:training-inference}
\noindent
Training proceeds in two stages, pre-training on the synthetic SSM cases and
fine-tuning on the clinical ED and ES cases, with the same architecture and the
same loss throughout. This section describes the two stages and the loss, then
how the trained fields are turned into meshes and an offset map, and finally
the two alternative reconstructions against which the model is compared.

\subsection{Two-stage training and loss}
\label{subsec:training-data-mix}
\noindent
In the first stage the model is trained only on the 1300 synthetic ED cases,
which give it a broad prior for LV shape under controlled conditions and
tolerate stronger augmentation of the sparse observations. In the second stage
it continues from the pre-trained weights on the ACDC and M\&Ms-2 cases of both
phases, adapting to clinical contour quality and to systolic anatomy. The
patient-level split and the augmentation rule of
\cref{subsec:shared-dataset-assembly} apply unchanged.

One phase-conditioned model is trained rather than one model per phase. The
reason follows from the data. Because the SSM provides ED geometry only, a
separate ES model could not have received equivalent pre-training, whereas a
shared model lets the ES cases benefit from the prior learned on the synthetic
ED corpus, with the phase label telling the encoder which configuration it is
seeing. The RBF-derived targets are constructed identically in both phases, so
no second asymmetry is introduced at fine-tuning. The cost is that phase
information is coarse, a binary label rather than a position in the cardiac
cycle, and that the prior remains biased towards diastolic shape.

Training adjusts the weights of the network so as to minimise a loss function,
a single number that measures how far the current predictions are from what
the targets require. Here that requirement has five parts, and each becomes
one term of the loss. The predicted surfaces should pass through the sampled
points of the target surfaces (surface term). The predicted fields should take
the stored signed-distance values at the query points throughout the volume,
not only on the surfaces (dense SDF term). Each field should behave like a
true distance, changing by one unit per unit travelled, which is expressed by
requiring its gradient to have unit length (Eikonal term). The fields should
not approach zero anywhere other than at the two surfaces, so values in free
space are pushed away from zero (off-surface term). And the offset $\delta$ at
the endocardial surface should equal the actual distance to the epicardial
target, so that the coupling learns the real wall separation rather than
merely keeping the fields ordered (wall term). The first four are standard in
implicit surface learning, following DeepSDF \parencite{deepsdf2019} and the
Eikonal regularisation of \textcite{gropp2020igr}; the fifth is specific to
the coupled output.

Formally, let $\mathcal{E}$ and $\mathcal{P}$ denote sampled endocardial and
epicardial surface points, $\mathcal{Q}$ the stored query points,
$\mathcal{R}$ the union of near-surface and free-space points used for the
gradient term, and $\mathcal{F}$ the free-space subset. At each
$q\in\mathcal{Q}$, $y_{\mathrm{endo}}(q)$ and $y_{\mathrm{epi}}(q)$ are the
stored signed-distance targets, and at each $x\in\mathcal{E}$ the wall target
$t(x)=-y_{\mathrm{epi}}(x)$ is the distance from that endocardial point to the
epicardial surface. Writing $\mathbb{E}$ for the mean over the indicated set
of points, the terms are
\begin{align}
	\mathcal{L}_{\mathrm{surf}} ={}&
	\mathbb{E}_{x\in\mathcal{E}}[|f_{\mathrm{endo}}(x)|]
	+ \mathbb{E}_{x\in\mathcal{P}}[|f_{\mathrm{epi}}(x)|], \notag \\
	\mathcal{L}_{\mathrm{SDF}} ={}&
	\mathbb{E}_{q\in\mathcal{Q}}[|f_{\mathrm{endo}}(q)-y_{\mathrm{endo}}(q)|]
	+ \mathbb{E}_{q\in\mathcal{Q}}[|f_{\mathrm{epi}}(q)-y_{\mathrm{epi}}(q)|], \notag \\
	\mathcal{L}_{\mathrm{eik}} ={}&
	\mathbb{E}_{x\in\mathcal{R}}[(\lVert\nabla_x f_{\mathrm{endo}}(x)\rVert_2-1)^2]
	+ \mathbb{E}_{x\in\mathcal{R}}[(\lVert\nabla_x f_{\mathrm{epi}}(x)\rVert_2-1)^2], \notag \\
	\mathcal{L}_{\mathrm{off}} ={}&
	\mathbb{E}_{x\in\mathcal{F}}[e^{-\alpha|f_{\mathrm{endo}}(x)|}]
	+ \mathbb{E}_{x\in\mathcal{F}}[e^{-\alpha|f_{\mathrm{epi}}(x)|}], \notag \\
	\mathcal{L}_{\mathrm{WT}} ={}&
	\mathbb{E}_{x\in\mathcal{E}}[\ell_{\beta_{\mathrm{WT}}}(\delta(x)-t(x))],
	\label{eq:training-loss-terms}
\end{align}
where $\alpha=5$ and $\ell_{\beta}$ is the Smooth-L1 penalty, which grows like
a squared error for small differences and like an absolute error for large
ones, so that a few badly wrong points do not dominate the term:
\begin{equation}
	\ell_{\beta}(a)=
	\begin{cases}
		a^2/(2\beta), & |a|<\beta, \\
		|a|-\beta/2, & |a|\geq\beta,
	\end{cases}
	\qquad \beta_{\mathrm{WT}}=0.02.
	\label{eq:wall-smooth-l1}
\end{equation}
All distances are expressed in the normalised coordinate system of each case,
and the complete objective is
\begin{align}
	\mathcal{L}_{\mathrm{total}} ={}&
	\lambda_{\mathrm{surf}}\mathcal{L}_{\mathrm{surf}} +
	\lambda_{\mathrm{SDF}}\mathcal{L}_{\mathrm{SDF}} +
	\lambda_{\mathrm{eik}}\mathcal{L}_{\mathrm{eik}} +
	\lambda_{\mathrm{off}}\mathcal{L}_{\mathrm{off}} +
	\lambda_{\mathrm{WT}}\mathcal{L}_{\mathrm{WT}} .
	\label{eq:training-objective}
\end{align}
where the weights $\lambda$ balance the five terms. The Eikonal term is
computed in full precision to keep its gradient calculation stable.
\Cref{tab:loss-terms} in \cref{app:supplementary-methods} summarises the role
of each term.

The training run reported here uses the AdamW optimiser, a standard
gradient-based method, with a learning rate (step size) of $2\times10^{-5}$
and weight decay $5\times10^{-4}$; the learning rate follows cosine annealing
with warm restarts, decreasing along a cosine curve and being reset at
intervals, with an initial restart period of 100 epochs (passes over the
training set). Pre-training lasted approximately 200 epochs; fine-tuning then
ran for 776 epochs, and the checkpoint at epoch 695 was selected by validation
loss. Training uses mixed precision (lower-precision arithmetic where it is
safe) and gradient clipping at 1.0 (a cap on the size of each update); each of
two GPUs processes a batch of 8 cases and gradients are accumulated over two
steps, for an effective batch size of 32. Weights are lightly clipped instead
of spectrally normalised, because the decoder is evaluated at many points per
step and GPU memory is the limiting resource. \Cref{tab:training-config}
collects these settings.

\begin{table}[H]
	\centering
	\caption[Model training configuration]{Training configuration used for the
	model.}
	\label{tab:training-config}
	\small
	\begin{tabular}{p{0.30\textwidth} p{0.48\textwidth}}
		\toprule
		Setting & Value \\
		\midrule
		Training strategy & SSM pre-training, then fine-tuning on real ED/ES \\
		Training mode & Combined ED and ES with phase conditioning \\
		Augmentation & Contour-input perturbations for synthetic ED and real ED/ES; targets unchanged \\
		Optimiser & AdamW \\
		Learning rate & $2\times10^{-5}$ \\
		Scheduler & Cosine annealing with warm restarts ($T_0=100$ epochs) \\
		Weight decay & $5\times10^{-4}$ \\
		Training duration & Approximately 200 pre-training epochs, then 776 fine-tuning epochs (best checkpoint: 695) \\
		Batch size & 8 per GPU $\times$ 2 GPUs, accumulated for 2 steps (effective batch size 32) \\
		Gradient clipping & 1.0 \\
		Precision & Automatic mixed precision \\
		Hardware & 2 GPUs in parallel \\
		Inference grid & $96^3$ \\
		\bottomrule
	\end{tabular}
\end{table}

\subsection{Inference, offset extraction, and comparators}
\label{subsec:inference-offset-extraction}
\noindent
At inference time, the encoder turns the contours into $z$ and $V$ once, and
the decoder is evaluated on a regular $96^3$ grid of points. Marching cubes
extracts the zero level sets of $f_{\mathrm{endo}}$ and $f_{\mathrm{epi}}$ from
this grid as the raw endocardial and epicardial meshes
\parencite{lorensen1987marchingcubes}; the mesh is thus a discretisation chosen
at inference time, not a property of the architecture, and could be extracted
at any resolution. For evaluation, faces of zero area are removed, only the
largest connected piece of each mesh is kept, holes are closed, and, where
needed, a mesh is rebuilt from a cleaned inside/outside voxel mask. This repair
produces usable geometry for surface and thickness measurement, but it also
means that the reported metrics describe a repaired mesh rather than the
unmodified decoder output.

Because the decoder predicts $\delta$ at every point, a raw offset is also
available at every endocardial vertex. In millimetres,
\begin{equation}
	o_{\mathrm{raw}}(x) = s\,\delta(x,z,v(x)),
	\label{eq:raw-analytic-offset}
\end{equation}
where $s$ is the scale that maps normalised coordinates back to millimetres.
For visualisation only, this offset is calibrated separately for each phase as
$o_{\mathrm{cal}} = \alpha\,o_{\mathrm{raw}}$, where $\alpha$ is the ratio of
the mean of a reference thickness, computed directly from the segmentation
voxels with the distance-transform method of \cref{subsec:thickness-methods},
to the mean raw offset. The calibration fixes the mean by
construction and is not a test of accuracy, and, for the reason given in
\cref{subsec:monotone-epi-decoder}, the offset itself is not a thickness
measurement; every reported thickness value comes from the protocol of
\cref{sec:wall-thickness-protocol}.

The model is compared with two reconstructions built from exactly the same
contour rings, chosen because they embody two different and well-understood
priors. The first is an RBF implicit fit obtained with the procedure of
\cref{subsec:real-ed-es-datasets}: it interpolates the rings through a smooth
field and belongs to the same family as the fine-tuning targets. The second is
a fit of the UK Digital Heart Project shape model to the rings. The mean shape
is first aligned to the rings by long axis, scale, and rotation about the long
axis, where the long axis is the line through the centres of the rings; the
fit then alternates three steps, matching each ring point to its nearest model
vertex, re-aligning the model to the rings by translation, rotation, and
scale, and re-estimating the first $25$ mode coefficients with the
coefficients constrained to three standard deviations. As the published model
provides ED
modes only, this comparator exists at ED alone. Both pass through the same
repair and smoothing as the model output.

Since no independent three-dimensional geometry exists, none of the three
reconstructions is treated as a reference. They are compared in two ways: by
the distance from the observed rings to each surface, which measures how
closely each method honours the data it was given, and by pairwise surface and
volume agreement, which measures how much the choice of prior changes the
result. Neither comparison isolates the contribution of individual model
components. The reconstructed meshes now become the input to the measurement
stage.

\section{Wall-Thickness Evaluation Protocol}
\label{sec:wall-thickness-protocol}
\noindent
The output of the previous section is a pair of surfaces; the clinical
quantity of interest is the distance between them. On a curved, tapering wall
that distance is not uniquely defined, because it depends on which point of
the epicardium is paired with a given point of the endocardium: the nearest
point, the point reached by travelling perpendicular to the surface, or the
point reached along a smooth path across the wall. Different pairing rules
give different values on the same geometry, and a rule that behaves well on
one geometry can be erratic on another. The protocol therefore keeps the two
sources of variation apart. Thickness is always measured on the reconstructed
endocardial and epicardial meshes; the segmentation volume provides only the
contour rings and the orientation of the AHA coordinate system. Four
established estimators with different pairing rules are run on the same
surfaces, and the one that proves most consistent in \cref{ch:results} is used
for the regional analysis. The protocol compares measurement methods and
constructs three-dimensional thickness maps; it does not assess diagnostic
performance for infarction, scarring, or hypertrophy.

\subsection{Selected thickness methods}
\label{subsec:thickness-methods}
\noindent
The four methods span the main families identified in
\cref{subsec:thickness-review}: distance fields computed on a voxel grid, rays
cast from one surface to the other, and methods based on partial differential
equations (PDEs), which define a smooth field across the wall and measure
along it (\cref{tab:wall-thickness-methods}). None is assumed in advance to be
the best. Each returns a local thickness in millimetres at the endocardial
vertices or at points projected onto the endocardium, and the resulting fields
are summarised by their mean, median, standard deviation, 5th and 95th
percentiles, fraction of valid values, and runtime, and shown as
three-dimensional colour maps and as AHA-17 bullseye plots, circular diagrams
with one cell per segment.

\begin{table}[H]
	\centering
	\caption[Wall-thickness methods selected for evaluation.]{Wall-thickness methods selected for evaluation on the
	model-generated LV geometry, with the Euclidean Distance Transform
	boundary sum kept as the simplest baseline.}
	\label{tab:wall-thickness-methods}
	\small
	\begin{tabular}{p{0.26\textwidth} p{0.18\textwidth} p{0.42\textwidth}}
		\toprule
		Method & Category & Main measurement principle \\
		\midrule
		Laplace field & PDE & Transmural field lines from a harmonic potential \\
		Yezzi--Prince & PDE & Eulerian transport equations along the transmural field \\
		SDF cone rays & Ray / SDF & Median of multiple cone rays around the surface normal \\
		EDT boundary sum & Volumetric & Sum of endocardial and epicardial distance transforms \\
		\bottomrule
	\end{tabular}
\end{table}

The Euclidean Distance Transform (EDT) boundary sum is the simplest. A
distance transform assigns to every voxel of the myocardium its straight-line
distance to the nearest point of a given boundary; the method computes one
such transform for each surface and adds them,
\begin{equation}
	t^{\mathrm{EDT}}(x) = D_{\mathrm{endo}}(x) + D_{\mathrm{epi}}(x),
	\label{eq:edt-thickness}
\end{equation}
where $D_{\mathrm{endo}}$ and $D_{\mathrm{epi}}$ are the distance transforms
of the endocardium and the epicardium. The sum is exact for two parallel flat
walls, but it is limited by voxel size and slightly overestimates thickness
where the wall is strongly curved, because the two nearest points do not lie
on a common path across the wall.

The Laplace field method replaces nearest points with a smooth path across the
wall. A scalar field $\psi$ is defined over the myocardium by solving Laplace's
equation,
\begin{equation}
	\nabla^2 \psi = 0,
	\qquad
	\psi|_{\mathrm{endo}} = 0,
	\qquad
	\psi|_{\mathrm{epi}} = 1.
	\label{eq:laplace-thickness}
\end{equation}
which makes $\psi$ rise as smoothly as possible from 0 on the endocardium to 1
on the epicardium, in the sense that its value at every point is the average
of the values around it; a field with this property is called harmonic, and
$\psi$ plays the role of a potential whose field lines run across the wall.
Thickness is then estimated from the magnitude of
the gradient of $\psi$, the rate at which the field changes across the wall,
which is small where the wall is thick and large where it is thin. The paths
that follow the gradient never cross one another, which is why the method,
introduced for the thickness of the brain cortex by \textcite{jones2000laplace},
is widely used for layered tissue; its weakness is that the numerical estimate
can degrade where the wall is thin or the gradient is small. The Yezzi--Prince
method \parencite{yezzi2003thickness} removes this weakness by computing,
along the same field, the distance travelled from each boundary with a pair of
transport equations solved on the voxel grid (an Eulerian formulation, meaning
that the quantities are computed at fixed grid points rather than by tracing
individual paths), in simplified form
\begin{equation}
	\nabla \psi \cdot \nabla u = -1,
	\qquad
	t = u + v,
	\label{eq:yezzi-prince-thickness}
\end{equation}
where $u$ and $v$ are the distances travelled from the two boundaries and
their sum is the thickness.

The cone-ray method pairs each endocardial point with the epicardial point
reached along its surface normal, but guards against a single badly oriented
ray. Following the shape diameter function of \textcite{shapira2008sdf}, $K$
rays are cast within a cone of half-angle $30^{\circ}$ around the normal, and
the median of the valid intersection distances is kept,
\begin{equation}
	t_i^{\mathrm{cone}} = \operatorname{median}_{k=1}^{K} d_{ik},
	\qquad K=7.
	\label{eq:cone-ray-thickness}
\end{equation}
The estimate keeps an approximately transmural interpretation, that is, it
measures across the wall rather than along it, while being more robust than a
single normal ray. \Cref{fig:wall-thickness-methods} contrasts the four pairing
rules on one wall cross-section.

\begin{figure}[H]
	\centering
	\includegraphics[width=\linewidth]{fig_wall_thickness_methods}
	\vspace{2mm}
	\caption[Comparison of wall-thickness estimators.]{Conceptual comparison of the four wall-thickness estimators on a
	single myocardial cross-section, with the endocardium in blue and the
	epicardium in orange. (a) The EDT boundary sum adds the distances from an
	interior point to each boundary. (b) The Laplace field method follows the
	curved paths of a smooth potential that rises from 0 on the endocardium to
	1 on the epicardium; the paths cross the level lines of the potential
	(dashed) at right angles, and thickness is estimated from the gradient
	magnitude. (c) The Yezzi--Prince method propagates two travel distances $u$
	and $v$ from opposite boundaries along that same curved field and sums
	them. (d) The SDF cone-ray method casts a fan of rays within a cone around
	the endocardial normal and takes the median intersection distance
	(highlighted), making it robust to any single ill-oriented ray.}
	\label{fig:wall-thickness-methods}
\end{figure}

What is specific to this work is not the estimators but how they are used: all
four are run on the same reconstructed surfaces under the same conditions, and
their agreement with one another, rather than with a reference that does not
exist, decides which one carries the regional analysis.

\subsection{Local three-dimensional wall thickness}
\label{subsec:local-3d-thickness}
\noindent
Before any regional summary, each estimator produces a thickness value at
every endocardial vertex of the repaired meshes of
\cref{subsec:inference-offset-extraction}, that is, a continuous map over the
whole reconstructed wall rather than a handful of measurements on selected
slices. This map is the primary output of the measurement stage. It shows the
gradual change from the thin apex to the thicker base and preserves local
thinning or thickening that averaging within an AHA segment would hide, and a
fixed colour scale makes the maps comparable across patients and phases
(\cref{fig:lv-wall-thickness-3d}).

\begin{figure}[H]
	\centering
	\includegraphics[width=0.9\linewidth]{fig_lv_wall_thickness_3d}
	\vspace{2mm}
	\caption[Local 3D wall-thickness map on the reconstructed LV.]{Local three-dimensional wall-thickness map on the reconstructed LV
	endocardial surface, shown from two opposing viewpoints. Each vertex is
	coloured by its local wall thickness on a fixed millimetre scale, so that the
	continuous spatial variation from the thin apex to the thicker base is
	visible directly on the geometry rather than only as regional averages.}
	\label{fig:lv-wall-thickness-3d}
\end{figure}

\subsection{Regional AHA-17 mapping}
\label{subsec:aha17-mapping}
\noindent
For clinical reporting, the local map of the selected estimator is summarised
in the AHA 17-segment model \parencite{cerqueira2002segmentation}. The model
was defined for two-dimensional imaging. The ventricle is divided along its
long axis into basal, mid-cavity, and apical thirds; each third is read on a
representative SAX slice and divided into angular sectors around the centre of
the cavity, six in the basal and mid-cavity thirds and four in the apical
third, numbered from the anterior wall with the septal segments facing the
right ventricle; and the tip of the apex forms the seventeenth segment. In
routine use, the thickness of each segment is then read on that slice at
points chosen by the observer.

Transferring this definition to a reconstructed surface requires only that
every vertex receive the two coordinates the definition relies on, a position
along the long axis and an angle around it; the segments then become patches
of the three-dimensional surface rather than readings on single slices. The
long axis is taken as the direction of the SAX stack. The base and apex are
the first and last slices that contain myocardium, the base being the end
whose cavity cross-section is larger, and every endocardial vertex is given a
normalised position between 0 at the base and 1 at the apex. Vertices in the
first third of this range form the basal ring, those in the second third the
mid-cavity ring, those between two thirds and 0.85 the apical ring, and those
beyond 0.85 the apical cap, which is segment 17. The angle of each vertex is
measured in the SAX plane around the centre of the endocardial vertices, with
its zero fixed by the right-ventricular label of the segmentation: the
direction from the LV centre to the centre of the right ventricle marks the
septum, the wall shared by the two ventricles, and the anterior reference lies
a quarter turn from it. The basal and mid-cavity rings are then cut into six
sectors of $60^{\circ}$ and the apical ring into four sectors of
$90^{\circ}$, giving segments 1--6, 7--12, and 13--16
(\cref{fig:aha17-cut}). A quality-assurance step visualises the detected
base, apex, and septal direction before the segments are formed. Unrolling
the rings onto concentric circles gives the familiar bullseye, with the basal
ring outermost and the apical cap at the centre; each cell of the bullseye now
summarises all vertices of one surface patch. Within each segment the vertex
values are combined with robust summaries, the mean and percentiles, so that
isolated outliers do not dominate the regional value. The aggregation does not
replace the per-vertex maps; it restates the same measurements in a form
clinicians are used to.

\begin{figure}[H]
	\centering
	\includegraphics[width=\linewidth]{fig_aha17_cut}
	\vspace{2mm}
	\caption[AHA-17 partition]{AHA-17 partitioning of the reconstructed LV
	surface and corresponding bullseye representation.}
	\label{fig:aha17-cut}
\end{figure}

This completes the pipeline of \cref{fig:methodology-pipeline}. Its outputs,
two reconstructed surfaces per case, four local thickness fields on them, and
AHA-17 summaries of the selected field, are the material of \cref{ch:results}.
There, agreement with the observed rings and with the RBF and SSM fits
addresses the first research question; the behaviour of the coupled surfaces
across ED and ES addresses the second; and the agreement among thickness
estimators and the regional comparisons address the third.
